\scanpage{56}
Je m'aperçois finalement que j'ai
pas mal \unsure{bousillé|bafouillé} pour arriver à trouver des
formulations simples --- car j'ai \uline{mélangé} des
propriétés purement ensemblistes de l'ens ordonné
$\Sigma=\Sigma^{<}(0)$, avec \unsure{\deleted{leur}|leur} interprétation semi-simpliciale,
et \deleted{\unsure{leur}}\inserted{les} aspects topologiques. Je \uline{reprends}

\deleted{$\Sigma$ ens.}
$\underset{\substack{\Vert\\\Sigma(0)}}{\Sigma}$ ens ordonné, muni d'une topologie \deleted{séparée}\inserted{(pour}
Si $\Sigma_1$ est le graphe de la relation d'ordre, on suppose

\textcircled{a} $\Sigma(1)\xrightarrow{\uline{b}\unsure{=|:}\operatorname{pr}_2}\Sigma(0)$ est propre et loc. immersif bordant
\inserted{(propriétés topologiques)} \deleted{Cela implique}

\deleted{\textcircled{b}} \deleted{que l'applic\unsure{ation}}
Cela implique que l'application diagonale est
ouverte, donc si $\Sigma$ séparé, son image est une
partie à la fois ouverte et fermée --- et on
peut alors introduire $\unsure{\Sigma_F^{<}(1)|\Sigma^{<}(1)}=\Sigma(1)\setminus\text{diagonale}$.
Si $\Sigma$ est séparé, la propreté de $\Sigma(1)\unsure{-|\longrightarrow}\Sigma(0)$
implique aussi que $\Sigma_1$ est \uuline{fermé} dans \deleted{$\Sigma(0)\times\Sigma(0)$}\inserted{$\Sigma\times\Sigma$}.

\textcircled{b} Pour tt $x\in\Sigma$, l'ens \inserted{(fini)} $\Sigma_{\leq x}$ \inserted{(fibre de $x$ dans $\Sigma(1)$)} est \uline{tot. ordonné}
pour l'ordre induit.
\deleted{\textcircled{\unsure{c}}} Donc la donnée d'un drapeau \inserted{$x_{i_0}-x_{i_d}$} de longueur
$d$ dans $\Sigma_\leq$ équivaut à celle d'une partie de card. $d+1$
\mspage{55}
